\documentclass[11pt]{article}
\usepackage{notes}        % custom package (optional)

\title{Problem Set: Fourier Transforms with Potential Fields}
\author{Geophysics 3A}
\date{}

\begin{document}
\maketitle

\section*{Overview}

In this practical you will use Fourier transforms to analyze gravity and magnetic
data from South Australia.  You will explore how spectral methods relate wavelength
to source depth, how continuation filters modify potential-field data, how
isostatic compensation can provide a physically motivated regional field, and how
reduction to the pole corrects for the inclination of the Earth's magnetic field.

Complete the accompanying Jupyter Notebook and answer the interpretation questions
below.  Your responses should demonstrate physical understanding rather than simply
restating computational steps.

\section*{Learning Objectives}

By completing this assignment, you should be able to:
\begin{itemize}
  \item Interpret potential-field data in both spatial and spectral domains
  \item Explain upward and downward continuation in physical and mathematical terms
  \item Separate regional and residual fields and critically evaluate the choices involved
  \item Compare spectral and isostatic approaches to defining a regional gravity field
  \item Apply reduction to the pole and understand its assumptions and limitations
  \item Discuss how sediment thickness affects the interpretation of magnetic data
  \item Recognise the non-uniqueness inherent in potential-field interpretation
\end{itemize}

% ─────────────────────────────────────────────────────────────────────────────
\section{One-Dimensional Profile Continuation}
% ─────────────────────────────────────────────────────────────────────────────

Gravity and magnetic profiles extracted from the gridded data are continued
upward and downward using a spectral filter.  The continuation filter is:
\[
  H(k,\, z) = e^{-2\pi |k| z}
\]
where $k$ is wavenumber (cycles~km$^{-1}$) and $z > 0$ is upward continuation
height in km.

\subsection*{Question 1 — Upward continuation as a low-pass filter}

As you increase the continuation height, which features disappear first ---
broad, regional anomalies or narrow, localised ones?
What does this tell you about the relationship between anomaly wavelength and
source depth?

\subsection*{Question 2 — Residual anomalies}

What geological features might the residual profile ($\text{observed} - \text{regional}$)
represent?  Are the residuals in the gravity and magnetic profiles spatially
correlated along your chosen profile?  What could explain agreement or
disagreement between the two fields?

\subsection*{Question 3 — Choosing a continuation height}

There is no single ``correct'' continuation height for separating regional from
residual.
\begin{enumerate}[label=\alph*.]
  \item What criteria might you use to make a defensible choice of continuation height?
  \item How sensitive is your geological interpretation to that choice?
\end{enumerate}

\subsection*{Question 4 — Instability of downward continuation}

Move the continuation-height slider into negative values (downward continuation).
\begin{enumerate}[label=\alph*.]
  \item At approximately what depth does the downward-continued field become visibly unstable?
  \item Does instability set in at the same depth for gravity and magnetics, or does one break down sooner?
  \item Why might they differ?
\end{enumerate}

\subsection*{Question 5 — Physical meaning of the continued field}

Upward continuation produces the field that \emph{would} be observed at a higher
elevation.  Does the continued profile look like what you would expect from a
smoother, more distant source distribution?  Are there any features that seem
physically unreasonable, and if so, what might cause them?

% ─────────────────────────────────────────────────────────────────────────────
\section{Two-Dimensional Gravity Continuation and Isostatic Correction}
% ─────────────────────────────────────────────────────────────────────────────

The 2-D continuation filter extends the 1-D result to two spatial frequencies:
\[
  H(k_x, k_y,\, z) = e^{-2\pi \sqrt{k_x^2 + k_y^2}\, z}
\]
An isostatic correction based on Airy compensation provides an alternative,
geologically motivated regional field.  Under Airy isostasy, topographic load
$h$ is compensated by a crustal root of thickness
$t_r = \rho_c\,h\,/\,(\rho_m - \rho_c)$, producing a gravity effect:
\[
  g_\mathrm{iso}(x,y) = 2\pi G\,\rho_c\, h(x,y)
\]

\subsection*{Question 6 — Spectral regional vs isostatic correction}

Compare the spectral regional field (at your chosen continuation height) with
the isostatic correction map.
\begin{enumerate}[label=\alph*.]
  \item How well do they agree spatially?
  \item Which areas show the largest differences, and why?
\end{enumerate}

\subsection*{Question 7 — Residual comparison}

At what continuation height does the spectral residual most closely resemble the
isostatic residual?  What does this imply about the effective depth of isostatic
compensation in South Australia?

\subsection*{Question 8 — Anomalies not explained by isostasy}

What geological features might explain anomalies that are present in the spectral
residual but absent (or reduced) in the isostatic residual?

\subsection*{Question 9 — Limitations of the isostatic model}

What assumptions in the Airy isostasy model most strongly limit its accuracy in
this region?  Consider both the model geometry and the physical properties used.

% ─────────────────────────────────────────────────────────────────────────────
\section{Reduction to the Pole}
% ─────────────────────────────────────────────────────────────────────────────

Reduction to the pole (RTP) applies a phase correction that transforms total
magnetic intensity (TMI) data as though the ambient field and induced
magnetisation were both vertical.  The filter is:
\[
  H_\mathrm{RTP}(k_x, k_y) =
    \frac{\bar{\Theta}^2}{\left(|\Theta|^2 + \varepsilon\right)^2},
  \quad
  \Theta = \sin I +
    i\cos I\!\left(\frac{k_x \cos D + k_y \sin D}{|\mathbf{k}|}\right)
\]
where $I$ and $D$ are the geomagnetic inclination and declination.

\subsection*{Question 10 — Effect of RTP on anomaly positions}

How does reduction to the pole change the positions and shapes of the major
magnetic anomalies in South Australia?  Are anomalies shifted northward or
southward, and by roughly how much?

\subsection*{Question 11 — Stability at low inclinations}

Experiment with the inclination slider.
\begin{enumerate}[label=\alph*.]
  \item At what inclination do artefacts (ringing, extreme amplitudes) begin to
        appear?
  \item What does this tell you about the reliability of RTP near the magnetic
        equator?
\end{enumerate}

\subsection*{Question 12 — Difference map}

Examine the difference map (RTP $-$ TMI).
\begin{enumerate}[label=\alph*.]
  \item How does it relate to the asymmetry of individual anomalies in the TMI?
  \item What geological information might be encoded in that difference?
\end{enumerate}

\subsection*{Question 13 — Remanent magnetisation}

The RTP filter assumes purely \emph{induced} magnetisation parallel to the
ambient field.  How would remanent magnetisation with a different direction affect
the result?  Can you identify any anomalies in the dataset that might carry
significant remanence?

% ─────────────────────────────────────────────────────────────────────────────
\section{Sediment-Thickness Correction for RTP [\emph{Optional}]}
% ─────────────────────────────────────────────────────────────────────────────

Thick, non-magnetic sedimentary cover attenuates basement magnetic signals.
A spectral correction attempts to undo this attenuation by multiplying the
spectrum by $e^{+2\pi|\mathbf{k}|\,z_s}$, where $z_s$ is the sediment thickness.
A Gaussian taper $e^{-2\pi^2|\mathbf{k}|^2\lambda^2}$ damps the resulting
noise amplification.

\subsection*{Question 14 — Where the correction matters most}

Where does the sediment correction make the largest difference to the RTP result?
Does this correspond to the areas of thickest sediment shown in the data?

\subsection*{Question 15 — Role of the damping parameter}

What happens when you increase the damping length $\lambda$?  What signal is
being suppressed, and is this geologically sensible?

\subsection*{Question 16 — Validity of a uniform correction}

The correction here uses the mean sediment thickness across the study area.
\begin{enumerate}[label=\alph*.]
  \item In which parts of the study area do you expect this uniform approximation
        to be most and least valid?
  \item What additional information would you need to make the correction
        spatially variable in a rigorous way?
\end{enumerate}

\subsection*{Question 17 — Analogous gravity correction}

Could a similar sediment-correction approach be applied to the gravity field?
How would it differ from the Airy isostatic correction computed in Section~2?

% ─────────────────────────────────────────────────────────────────────────────
\section{Synthesis}
% ─────────────────────────────────────────────────────────────────────────────

\subsection*{Question 18 — Wavelength and depth}

How do Fourier-domain methods relate signal wavelength to source depth?  State
the approximate relationship, identify its key assumptions, and describe one
situation in which it would give a misleading depth estimate.

\subsection*{Question 19 — Spectral vs isostatic regional field}

What are the advantages and disadvantages of using spectral continuation versus
isostatic correction to separate regional from residual gravity?  Under what
geological conditions would you prefer one over the other?

\subsection*{Question 20 — Non-uniqueness}

All potential-field methods suffer from the fundamental non-uniqueness of the
inverse problem.  Give one concrete example from this practical where two
different assumptions produced similar-looking results.  What additional data or
constraints could help discriminate between them?

% ─────────────────────────────────────────────────────────────────────────────
\section*{Submission Instructions}
% ─────────────────────────────────────────────────────────────────────────────

\begin{itemize}
  \item Provide written answers to the questions above, either within the notebook or as a separate document as specified by your instructor.
  \item Include screenshots/figures from your Jupyter Notebook as necessary to answer the questions.
  \item Figures should be clearly labelled and referred to in your written answers where appropriate.
\end{itemize}

\section*{Assessment Criteria}

Your work will be assessed on the basis of:
\begin{itemize}
  \item Physical understanding of Fourier methods and potential fields
  \item Correct and insightful interpretation of results
  \item Clarity and quality of written explanations
  \item Appropriate use of figures to support interpretation
\end{itemize}

\vspace{1em}
\noindent\textbf{End of Assignment}

\end{document}
